\section{Related Work}
\label{sec:related}

\paragraph{Vision-language-action policies.} A growing family of VLA models maps language and visual observations to robot actions, typically by fine-tuning a vision-language backbone with an action head: RT-2 \citep{brohan2023rt2}, OpenVLA \citep{kim2024openvla}, $\pi_0$ \citep{black2024pi0}, and SmolVLA \citep{shukor2025smolvla}, the compact flow-matching policy we study. These models are usually evaluated by aggregate task success on benchmarks such as LIBERO \citep{liu2023libero}. Our work is orthogonal to model design: we take a published policy and recipe as given and ask what its failures \emph{look like} geometrically, rather than proposing a new architecture or training objective.

\paragraph{Displacement robustness and the memory trap.} Closest to us, \citet{afi2025memorytrap} introduce a displacement protocol and show that VLA policies suffer a ``memory trap'': they fail when a manipulated object is moved from its training configuration. \citet{liberopro2025} similarly probe robustness of LIBERO-trained policies under controlled perturbations. We treat the displacement phenomenon and its protocol as established by this line of work---AFI already shows qualitatively that the hand is driven toward the original location. Our contribution is not to re-discover that failure, but (i) to \emph{quantify} its endpoint geometry with a displacement-relative reach field on a demonstrably competent policy, and (ii) to run a \emph{counterfactual AFI does not}: holding the physical object fixed while changing only its rendered position, testing whether the local rendered object position is sufficient to steer the reach. Where prior work observes the failure under physical displacement, we causally dissociate the rendered object position from the reach endpoint---a distinction the qualitative observation ``ignoring new spatial cues'' leaves open.

\paragraph{Action representation and equivariance.} Our intervention design exploits the translation behavior of relative (delta) action policies, for which observation-frame translations compose with the action map in a known way \citep{wang2025equivariant}. We use this equivariance only as an off-the-shelf property to justify leaving the wrist channel untouched in early intervention variants; we do not claim it as a contribution. More broadly, action-diffusion and flow-matching policies \citep{chi2023diffusion,black2024pi0,shukor2025smolvla} produce action \emph{chunks} rather than single steps, which shapes how we define and measure the reach endpoint.

\paragraph{Diagnostic and causal evaluation of policies.} Beyond aggregate success, a body of work probes \emph{why} learned controllers behave as they do through targeted perturbations and interventions. Our renderer-local object-paste probe is in this spirit: rather than translating the whole observation frame---which we find to be out-of-distribution for the policy and therefore confounded---we move only the target object's rendered appearance, so the intervention isolates the agentview object channel without changing the rest of the scene. To our knowledge, applying such a local, content-preserving visual intervention to dissociate the reach endpoint from object position in a VLA policy is new; we position it as a diagnostic tool rather than a claim about all policies.
